\section{HCA-DSE}
\label{sec:method}

\subsection{Partial architectures}
The exploration traverses a prefix tree over the decision vector \eqref{eq:decision}
under a fixed variable order $v_1,\dots,v_k$.

\begin{definition}[Partial architecture]
A partial architecture of depth $d$ is a prefix of assignments
$p=(a_1,\dots,a_d)$, $a_i\in D_i$. Its completions are
$\C(p)=\{x\in\X: x_i=a_i,\ i\le d\}$, with $|\C(p)|=\prod_{i>d}|D_i|$.
\end{definition}

Let $\sigma:\X\to\{0,1\}$ be the well-formedness predicate of the template: the number of
gateways does not exceed the number of edge nodes, the replication factor does not exceed
the number of edge nodes, and replication is undefined for policies that do not place the
ingest stage on the edge tier. The working space is $\X_\sigma=\{x:\sigma(x)=1\}$ and all
statements below refer to $\F_\sigma=\F\cap\X_\sigma$; the exhaustive baseline against
which we compare uses the same predicate. The Pareto front is
\begin{equation}
\PS=\{F(x)\ :\ x\in\F_\sigma,\ \nexists\,y\in\F_\sigma:\ F(y)\preceq F(x)\},
\end{equation}
where $u\preceq v$ denotes $u_i\le v_i$ for all $i$ and $u_j<v_j$ for at least one $j$.

\subsection{Admissible bounds}

\begin{definition}[Admissibility]
$\LB_i$ is admissible for objective $f_i$ if
$\LB_i(p)\le f_i(x)$ for every partial architecture $p$ and every $x\in\C(p)$.
\end{definition}

The partial-cost construction discussed by Neubauer et al. assumes that objectives are
\emph{assignment monotonous}, i.e.\ that every assignment can only increase the objective,
so that the partial cost is itself a bound~\cite{neubauer2018exact}. The waiting term of
\eqref{eq:latency} violates this assumption: it depends on $\rho_t$, which by
\eqref{eq:rho} is jointly determined by $n_t$, the class of the tier, $r$ and $\pi$, and
which \emph{decreases} when nodes are added. We therefore construct the bounds in two
steps.

\begin{lemma}[Term-wise relaxation]
\label{lem:relax}
Let $f_i(x)=\sum_{t\in T_i}\varphi_t(x)$ with $\varphi_t\ge0$, and let
\begin{equation*}
\underline{\varphi}_t(p)=\min_{x\in\C(p)}\varphi_t(x).
\end{equation*}
Then $\LB_i(p)=\sum_{t\in T_i}\underline{\varphi}_t(p)$ is admissible.
\end{lemma}
\begin{proof}
For every $x\in\C(p)$ and every $t$, $\underline{\varphi}_t(p)\le\varphi_t(x)$ by
definition of the minimum. Summing over $t\in T_i$ yields $\LB_i(p)\le f_i(x)$.
\end{proof}

The minima are taken \emph{independently per term}, so the combination of values that
realises them need not correspond to any single completion; the bound is therefore
optimistic but valid. The implementation uses domain extrema for service, communication,
static energy and cost, and minima of class-specific power-to-capacity ratios
for dynamic energy. Categorical policies are enumerated over their small
reachable set; they have no meaningful monotone endpoint order. Independent
extrema can relax a term further than its exact minimum, which remains safe.
The queueing contribution is bounded by zero until its determinants are fixed.
Domain aggregates are precomputed and repeated class-dependent terms cached;
cache misses can still require scans of reachable classes or policies.

\begin{lemma}[Exactness on determined terms]
\label{lem:exact}
Let $\mathrm{supp}(\varphi_t)$ be the set of variables $\varphi_t$ depends on. If
$\mathrm{supp}(\varphi_t)\subseteq\{v_1,\dots,v_d\}$ then $\varphi_t$ is constant on
$\C(p)$, and substituting its exact value for $\underline{\varphi}_t(p)$ preserves
admissibility while tightening the bound.
\end{lemma}
\begin{proof}
Constancy is immediate, and a constant equals its own minimum over $\C(p)$, so
Lemma~\ref{lem:relax} applies unchanged.
\end{proof}

\begin{remark}[Why the variable order matters]
\label{rem:order}
Once all determinants of a tier's waiting contribution are assigned,
that contribution is constant on $\C(p)$ and can enter the bound exactly by
Lemma~\ref{lem:exact}; before then, zero is admissible for stable completions.
If a determined used tier has $\rho_t\ge1$, its latency is infinite.
The default order O1 is $(n_e,\kappa_e,n_g,\kappa_g,\ell,\kappa_c,r,\pi)$.
Policy is assigned last, so its exact queueing contribution becomes available only at
depth eight. We separately test O2,
$(n_e,\kappa_e,n_g,\kappa_g,r,\pi,\ell,\kappa_c)$, which can fix edge and gateway
queueing determinants on incomplete designs. Section~\ref{sec:order} separates the
benefit of this term from the effect of changing traversal order.
\end{remark}

\subsection{Pruning rules}

\begin{proposition}[Safe feasibility pruning]
\label{prop:feas}
Let $\LB_{g_j}$ be admissible for constraint function $g_j$. If
$\LB_{g_j}(p)>g_j^{\max}$ then $\C(p)\cap\F=\varnothing$.
\end{proposition}
\begin{proof}
For any $x\in\C(p)$, $g_j(x)\ge \LB_{g_j}(p)>g_j^{\max}$, hence $x\notin\F$.
\end{proof}

\begin{proposition}[Safe dominance pruning]
\label{prop:dom}
Let $y\in\F_\sigma$ be a fully evaluated feasible architecture with
$F(y)\preceq \LB(p)=\big(\LB_1(p),\dots,\LB_n(p)\big)$. Then $F(x)\notin\PS$ for every
$x\in\C(p)$.
\end{proposition}
\begin{proof}
Let $x\in\C(p)$. By admissibility $f_i(x)\ge \LB_i(p)\ge F_i(y)$ for all $i$, and by
assumption there is a $j$ with $F_j(y)<\LB_j(p)\le f_j(x)$. Hence $y$ strictly dominates
$x$ and $F(x)$ is not on the Pareto front.
\end{proof}

\begin{proposition}[Closure of the well-formedness predicate]
\label{prop:struct}
Each invariant of $\sigma$ depends only on assigned variables and, under the chosen
order, all variables of an invariant are assigned before it is first checked. Since
assignments are never revised, $\neg\sigma(p)$ implies $\sigma(x)=0$ for every
$x\in\C(p)$, i.e.\ $\C(p)\cap\X_\sigma=\varnothing$.
\end{proposition}

\begin{remark}
Structural invariants define what counts as an architecture rather than following from
the constraints of the problem: the evaluator can formally compute objectives for, say, a
replicated deployment under a policy that does not place the ingest stage on the edge
tier, although such a configuration is meaningless. Structural pruning is therefore part
of the definition of the search space, and the exhaustive baseline applies the same
predicate. This distinction was found by the test described in
Section~\ref{sec:methodology}, which initially --- and incorrectly --- verified structural
pruning as a feasibility argument.
\end{remark}

\subsection{Algorithm}
Algorithm~\ref{alg:hca} traverses the prefix tree depth-first. At every node the three
rules are applied in increasing order of cost: structural validity, then feasibility
bounds, then dominance against an archive of non-dominated evaluated points. A leaf that
survives is evaluated in full; feasible leaves update the archive. Fig.~\ref{fig:tree}
illustrates the two bound-based rules.

\begin{algorithm}[t]
\caption{HCA-DSE}
\label{alg:hca}
\begin{algorithmic}[1]
\Require design space $\X$ with order $v_1..v_k$, workload $w$, evaluator $F$
\Ensure Pareto front $\mathcal{P}$
\State $A\gets\varnothing$ \Comment{archive of non-dominated evaluated vectors}
\State $S\gets\varnothing$ \Comment{evaluated feasible architectures}
\Procedure{Explore}{$p$, $d$}
  \If{$d=k$}
    \State $(F(p),\textit{feasible})\gets\textsc{Evaluate}(p,w)$
      \Comment{the expensive step}
    \If{\textit{feasible}} $S\gets S\cup\{p\}$; \textsc{ArchiveInsert}$(A,F(p))$ \EndIf
    \State \Return
  \EndIf
  \ForAll{$a\in D_{d+1}$}
    \State $q\gets p\cup\{v_{d+1}\!\mapsto\!a\}$
    \If{$\neg\sigma(q)$} \textbf{continue} \Comment{Prop.~\ref{prop:struct}} \EndIf
    \State $\LB\gets\textsc{Bounds}(q)$ \Comment{Lemmas~\ref{lem:relax},~\ref{lem:exact}}
    \If{$\exists j:\LB_{g_j}(q)>g_j^{\max}$} \textbf{continue}
      \Comment{Prop.~\ref{prop:feas}} \EndIf
    \If{$\exists f\in A: f\preceq \LB$} \textbf{continue}
      \Comment{Prop.~\ref{prop:dom}} \EndIf
    \State \textsc{Explore}$(q,d+1)$
  \EndFor
\EndProcedure
\State \textsc{Explore}$(\varnothing,0)$
\State \Return non-dominated subset of $\{F(x):x\in S\}$
\end{algorithmic}
\end{algorithm}

\begin{figure}[t]
\centering
\begin{tikzpicture}[font=\scriptsize, level distance=9mm,
  every node/.style={inner sep=1.5pt},
  n/.style={circle, draw, minimum size=3.4mm, inner sep=0pt},
  pr/.style={rectangle, draw=red!70!black, fill=red!8, rounded corners=1pt},
  ev/.style={rectangle, draw=edgec, fill=edgec!12, rounded corners=1pt},
  level 1/.style={sibling distance=30mm},
  level 2/.style={sibling distance=15mm},
  level 3/.style={sibling distance=13mm}]
\node[n] {}
  child {node[n] {$n_e{=}2$}
    child {node[pr] {$\LB_L>D$}
      edge from parent node[left, xshift=-1mm] {edge-low}}
    child {node[n] {}
      child {node[ev] {evaluate} edge from parent node[right, xshift=1mm] {$r{=}1$}}
      child {node[pr] {$F(y)\preceq \LB$} edge from parent node[right, xshift=2mm]
             {$r{=}3$}}
      edge from parent node[right, xshift=1mm] {edge-high}}}
  child {node[n] {$n_e{=}8$}
    child {node[pr] {$\LB_C>C^{\max}$} edge from parent node[left, xshift=-1mm]
           {fast link}}
    child {node[ev] {evaluate} edge from parent node[right, xshift=1mm] {std link}}};
\end{tikzpicture}
\caption{Pruning on the prefix tree. Sub-trees are discarded when the optimistic bound
already violates a constraint (Proposition~\ref{prop:feas}) or is dominated by an
evaluated feasible design (Proposition~\ref{prop:dom}); only surviving leaves reach the
expensive evaluator.}
\label{fig:tree}
\end{figure}

\subsection{Correctness}

\begin{theorem}[Pareto-front preservation]
\label{thm:exact}
If all bounds are admissible, pruning is performed only by
Propositions~\ref{prop:feas}--\ref{prop:struct}, and the archive contains only feasible
fully evaluated architectures, then the set of non-dominated vectors returned by
Algorithm~\ref{alg:hca} equals the true Pareto front:
$\mathcal{P}^{\mathrm{HCA}}=\PS$.
\end{theorem}
\begin{proof}
$(\subseteq)$ Every returned vector is feasible. If a returned vector
were dominated by any feasible vector, finiteness of $\X$ implies that it is
dominated by a Pareto-optimal vector. The following argument shows that every
Pareto-optimal vector is evaluated; it would therefore remove the dominated
vector during final filtering, a contradiction.

$(\supseteq)$ Let $u\in\PS$ and let $x\in\F_\sigma$ with $F(x)=u$. The traversal is a
prefix tree, so $x$ is reached through the unique chain
$p_0\subset p_1\subset\dots\subset p_k=x$. Suppose some $p_d$ were pruned. By
Proposition~\ref{prop:struct} this implies $\sigma(x)=0$ and by
Proposition~\ref{prop:feas} $x\notin\F$, both contradicting $x\in\F_\sigma$; by
Proposition~\ref{prop:dom} it implies the existence of a feasible $y$ strictly dominating
$x$, contradicting $u\in\PS$. Hence no $p_d$ is pruned, $x$ is evaluated, and $u$ survives
the final filtering.
\end{proof}

\begin{remark}[Duplicates]
The theorem is stated in objective space. The strict dominance rule does not prune a design merely because its
objective vector equals an archived vector. A separate duplicate-removal step may
retain one representative per vector without changing the front. All experiments therefore compare sets of objective
vectors.
\end{remark}

\begin{remark}[Order independence]
Proposition~\ref{prop:dom} holds for any archive of evaluated feasible architectures, so
the traversal order affects pruning efficiency but never correctness.
\end{remark}

\subsection{Complexity}
Let $C_E$ be the cost of a full evaluation, $C_P$ the cost of a bound computation,
$N_B$ the number of checked prefixes, $N_E$ the number of full evaluations and $|A|$
the maximum archive size. The search performs
\begin{equation}
T_{\mathrm{HCA}}=O\big(N_B(C_P+n|A|)+N_E C_E\big).
\label{eq:complexity}
\end{equation}
For measured cost projections we use $T_H=O_H+N_E C_E$ and
$T_X=O_X+N_X C_E$, where $N_X$ counts structurally valid designs evaluated by
exhaustive search, and $O_H,O_X$ are the respective measured non-evaluation overheads.
Provided $N_X>N_E$, the break-even condition is
\begin{equation}
C_E > \frac{O_H-O_X}{N_X-N_E}.
\label{eq:breakeven}
\end{equation}
A negative threshold means HCA-DSE is already favoured at zero evaluator cost in this
linear model. Using $|\X|$ in place of $N_X$ would incorrectly charge exhaustive search
for structurally rejected designs. Expanded partial states $N$, reported separately,
exclude branches rejected before recursion; $N_B$ counts their bound checks as well.
The Cartesian product is never materialised by HCA-DSE.
